\documentclass[11pt]{article}

\usepackage[utf8]{inputenc}
\usepackage[T1]{fontenc}
\usepackage{lmodern}
\usepackage{amsmath,amssymb}
\usepackage{geometry}
\usepackage{microtype}
\usepackage{enumitem}
\geometry{margin=1in}

\newcommand{\R}{\mathbb{R}}
\newcommand{\F}{\mathcal{F}}
\newcommand{\eps}{\varepsilon}
\newcommand{\meas}[1]{\lvert #1\rvert}

\title{Meeting Summary: Almost Equicontinuity and Truncated Translation Conditions}
\author{Hikmatullo Ismatov}
\date{}

\begin{document}

\maketitle

\section*{Main Question}

The goal is to understand when almost equicontinuity implies the truncated
translation condition
\[
    \lim_{|y|\to0}\sup_{f\in\F}
    \int_{\R^n}\min\bigl(|f(x+y)-f(x)|,1\bigr)\,dx=0,
\]
and conversely, when this translation condition forces almost equicontinuity.
The truncation is important because it allows estimates on exceptional sets
without any pointwise boundedness or uniform $L^p$ assumption.

\section*{Definitions Used}

A family $\F$ of real-valued measurable functions on $\R^n$ is almost
equicontinuous if, for every $\eps>0$, there exists $\delta>0$ such that for
each $f\in\F$ there is a measurable set $B_f\subset\R^n$ with
$\meas{B_f}<\eps$ and
\[
    x,z\in\R^n\setminus B_f,\quad |x-z|<\delta
    \quad\Longrightarrow\quad
    |f(x)-f(z)|<\eps .
\]

The family $\F$ has uniformly tight supports if
\[
    \lim_{R\to\infty}\sup_{f\in\F}
    \meas{\{x\in\R^n: |x|>R,\ f(x)\ne0\}}=0.
\]

\section*{Results Proved}

\begin{enumerate}[label=\textbf{R\arabic*.},leftmargin=*]
    \item \textbf{Finite-measure forward implication.}
    If every $f\in\F$ is supported in a fixed measurable set
    $E\subset\R^n$ with $\meas{E}<\infty$, and if $\F$ is almost
    equicontinuous, then $\F$ satisfies the truncated translation condition.

    \item \textbf{The finite-measure assumption is real.}
    On $E=\R$, the singleton family $\F=\{\sin x\}$ is almost
    equicontinuous, but for every small nonzero $y$,
    \[
        \int_{\R}\min\bigl(|\sin(x+y)-\sin x|,1\bigr)\,dx=\infty.
    \]
    Thus one cannot simply remove the finite-measure support assumption.

    \item \textbf{Forward implication under uniformly tight supports.}
    If $\F$ is almost equicontinuous and has uniformly tight supports, then
    $\F$ satisfies the truncated translation condition on all of $\R^n$.

    \item \textbf{Finite-measure converse on the domain.}
    Let $E\subset\R^n$ have finite measure, and let the functions be defined
    on $E$. If
    \[
        \lim_{|y|\to0}\sup_{f\in\F}
        \int_{E\cap(E-y)}
        \min\bigl(|f(x+y)-f(x)|,1\bigr)\,dx=0,
    \]
    then $\F$ is almost equicontinuous on $E$.

    \item \textbf{Local zero-extension converse.}
    If $E$ is arbitrary and the zero extensions satisfy the global truncated
    translation condition on $\R^n$, then the restrictions of $\F$ to every
    finite-measure set $K\subset E$ are almost equicontinuous on $K$.

    \item \textbf{Equivalence under uniformly tight supports.}
    If $\F$ has uniformly tight supports, then
    \[
        \F\hbox{ is almost equicontinuous}
        \quad\Longleftrightarrow\quad
        \F\hbox{ satisfies the truncated translation condition}.
    \]
\end{enumerate}

\section*{Proof Ideas}

\paragraph{Forward direction on finite-measure support.}
The integral is localized to $E\cup(E-y)$. This region is split into
\[
    E\triangle(E-y)
    \quad\hbox{and}\quad
    E\cap(E-y).
\]
The symmetric-difference part is small because finite-measure sets are
continuous under translation. On the intersection, remove
$B_f\cup(B_f-y)$. Outside this exceptional set, almost equicontinuity controls
the integrand; on the exceptional set, the truncation by $1$ controls the
integral by its measure.

\paragraph{Uniform tail version.}
Choose a large ball so that the support outside the ball is uniformly small.
Inside the ball, use the finite-measure argument. Outside the ball, the
truncation and the tail-support estimate control the contribution.

\paragraph{Finite-measure converse.}
The translation condition first implies that, for every fixed threshold
$a>0$, the bad set
\[
    \{x\in E\cap(E-y): |f(x+y)-f(x)|\ge a\}
\]
has measure tending to zero uniformly in $f$ as $y\to0$.

To prove almost equicontinuity by contradiction, one assumes that no small
exceptional set removes all nearby bad pairs. A shifted cube decomposition is
then used. On many cubes, one must remove a positive amount of measure before
the values of $f$ fit into an interval of length $\eps/2$. A separated-values
lemma converts this into a lower bound on the measure of bad pairs. Finally,
the change of variables $z=x+y$ converts bad pairs into an averaged lower
bound over translations, contradicting the translation condition.

\section*{Key Point to Explain Clearly}

The most delicate step is the passage
\[
    \hbox{many nearby bad pairs}
    \quad\Longrightarrow\quad
    \hbox{one small translation with many bad points}.
\]
This is done by introducing
\[
    P_s=
    \{(x,z)\in K\times K: |x-z|\le \sqrt n\,s,\,
    |f_s(x)-f_s(z)|\ge\eps/2\}
\]
and proving $\meas{P_s}\gtrsim s^n$. Writing $z=x+y$ gives
\[
    \int_{|y|\le\sqrt n\,s}
    \meas{\{x: |f_s(x+y)-f_s(x)|\ge\eps/2\}}\,dy
    \gtrsim s^n.
\]
Since the ball $\{|y|\le\sqrt n\,s\}$ also has measure comparable to $s^n$,
there is some $y_s\to0$ for which the bad set has measure bounded below by a
positive constant. This contradicts the translation condition.

\section*{Open Direction}

The remaining question is to identify the weakest global hypothesis replacing
uniformly tight supports. The local zero-extension converse shows that the
translation condition gives almost equicontinuity on every finite-measure
piece. What is not yet clear is whether these local exceptional sets can be
assembled into a single global exceptional set without a tightness assumption,
or whether there is a counterexample where oscillations escape to infinity.

\end{document}
